\section{Etale coverings}
\label{mumford-IV18}
\dcref{M:ch4:IV.18}

\begin{theorem}[Etale covers inherit the abelian structure]
\label{mumford-thm-etale-cover-abelian-structure}
Let $A$ be an abelian variety and let $u:Y\to A$ be a connected finite etale
cover with a chosen point above $0$.  There is a unique group law on $Y$ for
which that point is zero and $u$ is a homomorphism; $u$ is then a separable
isogeny.
\dcref{M:ch4:IV.18}
\uses{mumford-def-abelian-variety,mumford-prop-abelian-quotient}
\notready
\end{theorem}

\begin{lemma}[Graph construction for an etale cover]
\label{mumford-lem-etale-cover-graph}
Pull back the graph of multiplication on $A$ along $u\times u\times u$.
The component through the selected zero projects isomorphically to
$Y\times Y$; its third projection supplies the multiplication in
Theorem~$\ref{mumford-thm-etale-cover-abelian-structure}$.  The
fibre-irreducibility step is an external input.
\dcref{M:ch4:IV.18}
\uses{mumford-thm-etale-cover-abelian-structure}
\notready
\end{lemma}

\begin{definition}[Tate module]
\label{mumford-def-tate-module}
\dcref{M:ch4:IV.18}
Let $A$ be an abelian variety over an algebraically closed field and let
$\ell\ne\operatorname{char}(k)$.  Its $\ell$-adic Tate module is
\[
 T_\ell(A)=\varprojlim_n A[\ell^n](k),
\]
with transition maps multiplication by $\ell$.  It is a free $\Zell$-module
of rank $2\dim A$.
\uses{mumford-prop-multiplication-n,mumford-thm-abelian-smooth,
 mumford-thm-etale-cover-abelian-structure}
\notready
\end{definition}

\begin{proposition}[Etale coverings and torsion]
\label{mumford-prop-etale-coverings-av}
Every connected finite etale cover of $A$ is the quotient by a finite subgroup
of translations (equivalently, it is a separable isogeny after choosing a
point over the origin).  The prime-to-$p$ part of the
geometric fundamental group is
\[
 \pi_1(A)^{(p')}\simeq\prod_{\ell\ne p}\Zell^{\,2\dim A},
\]
and in characteristic zero this is the full profinite completion of
$H_1(A(\C),\Z)$.
\dcref{M:ch4:IV.18}
\uses{mumford-def-tate-module,mumford-prop-abelian-quotient}
\notready
\end{proposition}

\begin{proposition}[Galois action]
\label{mumford-prop-galois-tate-action}
\dcref{M:ch4:IV.18}
If $A$ is defined over a field $K$, the absolute Galois group $G_K$ acts
continuously on $T_\ell(A)$, giving a representation
\[
 \rho_{A,\ell}:G_K\to\Aut_{\Zell}(T_\ell(A)).
\]
Every homomorphism of abelian varieties over $K$ is $G_K$-equivariant on Tate
modules.
\uses{mumford-def-tate-module}
\notready
\end{proposition}

\section[Structure of Hom(X, dual X)]{Structure of homomorphisms into the dual}
\label{mumford-IV19}
\dcref{M:ch4:IV.19}

\begin{definition}[Hom groups up to isogeny]
\label{mumford-def-hom-up-to-isogeny}
\dcref{M:ch4:IV.19}
For abelian varieties $A,B$, set
\[
 \Hom^0(A,B)=\Hom(A,B)\otimes_\Z\Q,\qquad
 \End^0(A)=\End(A)\otimes_\Z\Q.
\]
Composition makes abelian varieties and these rational Hom groups into a
category in which every isogeny is invertible.
\uses{mumford-thm-isogeny-factorization}
\notready
\end{definition}

\begin{theorem}[Tate-module faithfulness]
\label{mumford-thm-hom-tate-injective}
For $\ell\ne\operatorname{char}(k)$, the natural map
\[
 \Hom(A,B)\otimes\Zell\longrightarrow
 \Hom_{\Zell}(T_\ell(A),T_\ell(B))
\]
is injective.  Consequently $\Hom(A,B)$ is torsion-free and finitely
generated, hence free abelian of finite rank.
\dcref{M:ch4:IV.19}
\uses{mumford-def-tate-module,mumford-def-hom-up-to-isogeny}
\notready
\end{theorem}

\begin{theorem}[Poincare complete reducibility]
\label{mumford-thm-poincare-complete-reducibility}
If $B\subset A$ is an abelian subvariety, there is an abelian subvariety
$C\subset A$ such that $B\cap C$ is finite and $B+C=A$.  Thus $A$ is isogenous
to $B\times C$, and every abelian variety is isogenous to a product of simple
ones.
\dcref{M:ch4:IV.19}
\uses{mumford-def-polarization,mumford-thm-ample-finite-kernel}
\notready
\end{theorem}

\begin{corollary}[Endomorphism algebra of a simple variety]
\label{mumford-cor-simple-endomorphism-algebra}
\dcref{M:ch4:IV.19}
If $A$ is simple, $\End^0(A)$ is a finite-dimensional division algebra over
$\Q$.  If its center is $K$, and its rank over $K$ is $d^2$, then
$d^2[K:\Q]$ divides $(2\dim A)^2$; in characteristic zero the complex
multiplication condition is the existence of a commutative subfield of degree
$2\dim A$.
\uses{mumford-thm-poincare-complete-reducibility,mumford-thm-hom-tate-injective}
\notready
\end{corollary}

\begin{remark}[Neron--Severi/Rosati package]
\label{mumford-iv20-ns-rosati-symmetric}
After tensoring with $\Q$, the Neron--Severi group identifies with the
Rosati-symmetric endomorphisms.
\dcref{M:ch4:IV.20}
\uses{mumford-iv20-riemann-form-homomorphism-criterion,mumford-prop-pic-ns-sequence}
\notready
\end{remark}

\section{Supporting dependency interfaces}
\label{mumford-IV-coverage-register}
The following nodes expose proof interfaces and examples used by the duality and
l-adic constructions.  Detailed proofs are treated as external inputs where
the available evidence is insufficient for reconstruction.

\begin{lemma}[Fibre irreducibility]
\label{mumford-iv18-fiber-irreducibility-lemma}
For a proper smooth morphism with a section, the fibre through the selected
point is irreducible; an idempotent in the completed local ring would otherwise
contradict locality.
\dcref{M:ch4:IV.18}
\uses{mumford-lem-etale-cover-graph}
\notready
\end{lemma}

\begin{remark}[Isogeny quasi-inverse]
\label{mumford-iv18-isogeny-quasi-inverse}
Every isogeny admits a quasi-inverse after tensoring Hom with $\Q$, with
integral composites equal to multiplication by a positive integer.
\dcref{M:ch4:IV.18}
\uses{mumford-thm-isogeny-factorization}
\notready
\end{remark}

\begin{definition}[Fundamental-group inverse system]
\label{mumford-iv18-fundamental-group-inverse-system}
The inverse limit of deck groups of pointed finite etale covers defines
$\pi_1(A)$.  Its prime-to-$p$ part is the product of the Tate modules, and in
characteristic zero it is the profinite completion of the analytic lattice.
\dcref{M:ch4:IV.18}
\uses{mumford-thm-etale-cover-abelian-structure,mumford-def-tate-module}
\notready
\end{definition}

\begin{definition}[Hom up to isogeny]
\label{mumford-iv19-hom-up-to-isogeny}
$\Hom^0(A,B)=\Hom(A,B)\otimes_\Z\Q$ and
$\End^0(A)=\End(A)\otimes_\Z\Q$ form the isogeny category.
\dcref{M:ch4:IV.19}
\uses{mumford-def-hom-up-to-isogeny,mumford-iv18-isogeny-quasi-inverse}
\notready
\end{definition}

\begin{proposition}[Complex Hom lattice]
\label{mumford-iv19-complex-hom-lattice}
For complex uniformizations, an algebraic homomorphism lifts to a complex-linear
map carrying one period lattice into the other; hence Hom is a finitely
generated lattice of rank at most $4\dim A\dim B$.
\dcref{M:ch4:IV.19}
\uses{mumford-prop-torus-homomorphisms,mumford-iv19-hom-up-to-isogeny}
\notready
\end{proposition}

\begin{definition}[The l-adic Hom representation]
\label{mumford-iv19-ladic-representation}
Restriction to torsion gives
\[
 \Hom(A,B)\longrightarrow\Hom_{\Zell}(T_\ell A,T_\ell B),
\]
compatible with composition and Galois action.
\dcref{M:ch4:IV.19}
\uses{mumford-def-tate-module,mumford-prop-galois-tate-action}
\notready
\end{definition}

\begin{theorem}[Endomorphism characteristic polynomial]
\label{mumford-iv19-end-characteristic-polynomial}
For $u\in\End(A)$ there is a monic $P_u\in\Z[t]$ of degree $2\dim A$,
independent of $\ell$, with $P_u(T_\ell u)=0$ and
$\deg(n-u)=P_u(n)$.
\dcref{M:ch4:IV.19}
\uses{mumford-thm-hom-tate-injective,mumford-iv19-ladic-representation}
\notready
\end{theorem}

\begin{theorem}[Simple endomorphism algebra]
\label{mumford-iv19-simple-endomorphism-algebra}
For a simple $A$, $\End^0(A)$ is semisimple (a division algebra in the simple
case), and its center/index satisfy the dimension divisibility constraints
coming from the characteristic polynomial.
\dcref{M:ch4:IV.19}
\uses{mumford-iv19-end-characteristic-polynomial,
 mumford-thm-poincare-complete-reducibility}
\notready
\end{theorem}

\begin{definition}[Tate pairing]
\label{mumford-iv20-tate-pairing}
The Poincare bundle supplies compatible perfect pairings
$A[\ell^n]\times\Adual[\ell^n]\to\mu_{\ell^n}$ and an inverse-limit pairing
$T_\ell A\times T_\ell\Adual\to\Zell(1)$.
\dcref{M:ch4:IV.20}
\uses{mumford-def-weil-pairing-ladic}
\notready
\end{definition}

\begin{definition}[Riemann form]
\label{mumford-iv20-riemann-form-definition}
The alternating form associated with a line bundle is the Tate-module form
obtained from the pairing and $T_\ell(\phi_L)$, after a choice of Tate-twist
generator when an additive value is desired.
\dcref{M:ch4:IV.20}
\uses{mumford-def-riemann-form-algebraic,mumford-iv20-tate-pairing}
\notready
\end{definition}

\begin{theorem}[Skew-symmetry criterion]
\label{mumford-iv20-riemann-form-skew}
The Tate-module form attached to $u:A\to\Adual$ is alternating exactly when
$u$ is represented by a line bundle; ampleness makes it a polarization form.
\dcref{M:ch4:IV.20}
\uses{mumford-iv20-riemann-form-definition,mumford-thm-riemann-form-hom}
\notready
\end{theorem}

\begin{theorem}[Riemann-form homomorphism criterion]
\label{mumford-iv20-riemann-form-homomorphism-criterion}
Skew forms on the Tate module correspond to homomorphisms into the dual, and
the image of $\NS(A)$ is the Rosati-symmetric part of $\End^0(A)$.
\dcref{M:ch4:IV.20}
\uses{mumford-iv20-riemann-form-skew,mumford-prop-pic-ns-sequence}
\notready
\end{theorem}

\begin{theorem}[Exterior-product generator]
\label{mumford-iv20-riemann-wedge-generator}
Exterior powers of a Riemann form generate the top alternating line and recover
intersection numbers and Euler characteristics.
\dcref{M:ch4:IV.20}
\uses{mumford3-riemann-roch,mumford-iv20-riemann-form-definition}
\notready
\end{theorem}

\begin{theorem}[Degree determinant formula]
\label{mumford-iv20-degree-determinant}
The degree of an endomorphism is the determinant of its action on a Tate module;
the same determinant description applies to the associated Riemann form.
\dcref{M:ch4:IV.20}
\uses{mumford-iv19-end-characteristic-polynomial,mumford-iv20-tate-pairing}
\notready
\end{theorem}

\begin{theorem}[Albert classification interface]
\label{mumford-iv21-albert-classification}
A simple rational endomorphism algebra with its Rosati-positive involution is
of one of Albert's four types, subject to the center and Schur-index bounds.
\dcref{M:ch4:IV.21}
\uses{mumford-thm-rosati-positive,mumford-iv19-simple-endomorphism-algebra}
\notready
\end{theorem}

\begin{theorem}[Weil eigenvalues and point counts]
\label{mumford-iv21-weil-rational-point-count}
For $A/\Fq$, Frobenius eigenvalues have complex absolute value $q^{1/2}$,
occur in reciprocal pairs, and
$\#A(\mathbb{F}_{q^n})=\prod_i(1-\omega_i^n)$.
\dcref{M:ch4:IV.21}
\uses{mumford-def-frobenius-polynomial,mumford-thm-frobenius-semisimple}
\notready
\end{theorem}

\begin{theorem}[Finite-field rational-point finiteness]
\label{mumford-iv21-lang-rational-point}
The finite-field point-count and Frobenius estimates supply the finiteness
inputs used in the endomorphism and torsion arguments.
\dcref{M:ch4:IV.21}
\uses{mumford-iv21-weil-rational-point-count}
\notready
\end{theorem}

\begin{theorem}[Torsion restriction]
\label{mumford-iv21-serre-torsion-restriction}
For a nonzero endomorphism of an abelian variety, restriction to sufficiently
large prime-to-characteristic torsion is nontrivial; this makes the Tate action
faithful.
\dcref{M:ch4:IV.21}
\uses{mumford-thm-hom-tate-injective}
\notready
\end{theorem}

\begin{corollary}[Endomorphism type restrictions]
\label{mumford-iv21-endomorphism-type-table}
The Albert type, center degree, and Schur index of a simple factor satisfy the
dimension restrictions imposed by its positive Rosati involution.
\dcref{M:ch4:IV.21}
\uses{mumford-iv21-albert-classification}
\notready
\end{corollary}

\begin{example}[CM type over $\C$]
\label{mumford-iv22-cm-complex-example}
A complex abelian variety with a commutative endomorphism field of degree twice
its dimension is of CM type; its Riemann form is constrained by the CM action.
\dcref{M:ch4:IV.22}
\uses{mumford-iv19-simple-endomorphism-algebra,mumford-thm-riemann-comparison}
\notready
\end{example}

\begin{example}[Elliptic endomorphisms in characteristic $p$]
\label{mumford-iv22-elliptic-weierstrass-example}
Elliptic curves have rational endomorphism algebra $\Q$, an imaginary
quadratic field, or a quaternion algebra in the supersingular case.
\dcref{M:ch4:IV.22}
\uses{mumford-ex-elliptic-endomorphisms}
\notready
\end{example}

\begin{theorem}[Deuring equivalences]
\label{mumford-iv22-deuring-equivalences}
For elliptic curves over finite fields, the endomorphism and Frobenius
polynomial criteria give the standard ordinary/supersingular isogeny
equivalences.
\dcref{M:ch4:IV.22}
\uses{mumford-iv21-weil-rational-point-count,mumford-thm-frobenius-isogeny-criterion}
\notready
\end{theorem}

\begin{theorem}[CM and finite-field isogeny]
\label{mumford-iv22-cm-finite-field-isogeny}
Equality of the Frobenius polynomial is equivalent to isogeny over a finite
field, including the CM and supersingular elliptic examples.
\dcref{M:ch4:IV.22}
\uses{mumford-thm-frobenius-isogeny-criterion,mumford-iv22-deuring-equivalences}
\notready
\end{theorem}

\begin{lemma}[Theta-group prime-order commutativity]
\label{mumford-iv23-theta-group-triviality}
For a finite theta group of prime exponent, the commutator pairing controls
whether a subgroup is central and whether the extension splits locally.
\dcref{M:ch4:IV.23}
\uses{mumford-def-theta-group,mumford-prop-theta-commutator}
\notready
\end{lemma}

\begin{theorem}[Line-bundle theta group]
\label{mumford-iv23-line-bundle-theta-group}
Automorphisms of a line bundle over translations form the central extension
$1\to\Gm\to\mathcal G(L)\to K(L)\to1$.
\dcref{M:ch4:IV.23}
\uses{mumford-def-theta-group}
\notready
\end{theorem}

\begin{definition}[Divisor description of a theta group]
\label{mumford-iv23-divisor-description}
Choosing a divisor representative for $L$ identifies theta-group elements with
translation/divisor pairs and makes the commutator explicit.
\dcref{M:ch4:IV.23}
\uses{mumford-iv23-line-bundle-theta-group}
\notready
\end{definition}

\begin{theorem}[Theta lifting and descent]
\label{mumford-iv23-isogeny-lifting-theorem}
Descent of $L$ along an isogeny is equivalent to lifting its kernel into the
theta group; the descended bundle has the expected pullback.
\dcref{M:ch4:IV.23}
\uses{mumford-iv23-line-bundle-theta-group,mumford-prop-abelian-quotient}
\notready
\end{theorem}

\begin{theorem}[Power criterion for line bundles]
\label{mumford-iv23-power-criterion}
For nondegenerate $L$, containment of $A[n]$ in $K(L)$ is equivalent to $L$
being an $n$th tensor power after the corresponding descent obstruction
vanishes.
\dcref{M:ch4:IV.23}
\uses{mumford-iv23-isogeny-lifting-theorem,mumford-prop-multiplication-n}
\notready
\end{theorem}

\begin{theorem}[Maximal isotropic theta subgroup]
\label{mumford-iv23-maximal-isotropic-theorem}
For a nondegenerate theta group, a maximal isotropic subgroup equals its
commutator orthogonal and has square-root order in $K(L)$; this controls
principal polarizations and the irreducible theta representation.
\dcref{M:ch4:IV.23}
\uses{mumford-prop-theta-commutator,mumford-thm-theta-representation}
\notready
\end{theorem}

\begin{definition}[Analytic theta extension]
\label{mumford-iv24-analytic-theta-extension}
For a complex torus and Hermitian form $H$, automorphisms of the factor of
automorphy form a central Heisenberg extension whose commutator is determined
by $\operatorname{Im}H$ on the lattice.
\dcref{M:ch4:IV.24}
\uses{mumford-thm-appell-humbert,mumford-def-theta-group}
\notready
\end{definition}

\begin{theorem}[Analytic/l-adic pairing comparison]
\label{mumford-iv24-analytic-ladic-comparison}
Under the comparison between the analytic lattice and the Tate module, the
analytic theta commutator agrees with the Poincare pairing up to the fixed sign
convention.
\dcref{M:ch4:IV.24}
\uses{mumford-thm-riemann-comparison,mumford-iv20-tate-pairing}
\notready
\end{theorem}

\begin{theorem}[Analytic Rosati positivity]
\label{mumford-iv24-analytic-rosati}
For a positive Hermitian form, the Rosati involution is the Hermitian adjoint;
its trace form is positive definite, giving the analytic proof of Rosati
positivity.
\dcref{M:ch4:IV.24}
\uses{mumford-thm-riemann-comparison,mumford-thm-rosati-positive}
\notready
\end{theorem}

\begin{remark}[Serre--Lang structure theorem]
\label{mumford-iv18-etale-covering-abelian-structure}
The Serre--Lang theorem asserts that every connected finite etale cover in the
etale-covering construction acquires the abelian-variety structure displayed
in Theorem~$\ref{mumford-thm-etale-cover-abelian-structure}$.
\dcref{M:ch4:IV.18}
\uses{mumford-thm-etale-cover-abelian-structure}
\notready
\end{remark}

\begin{remark}[Etale graph-pullback construction]
\label{mumford-iv18-etale-graph-pullback}
The pulled-back multiplication graph and its distinguished component implement
the group law on the cover.
\dcref{M:ch4:IV.18}
\uses{mumford-lem-etale-cover-graph}
\notready
\end{remark}

\begin{remark}[Complete reducibility package]
\label{mumford-iv19-poincare-complete-reducibility}
Every abelian subvariety has an isogeny complement, so an abelian variety is
isogenous to a product of simple factors.
\dcref{M:ch4:IV.19}
\uses{mumford-thm-poincare-complete-reducibility}
\notready
\end{remark}

\begin{remark}[Finite-generation and injectivity package]
\label{mumford-iv19-hom-finite-generation-injectivity}
Hom groups are torsion-free finitely generated and their l-adic realization is
faithful for every prime away from the characteristic.
\dcref{M:ch4:IV.19}
\uses{mumford-thm-hom-tate-injective,mumford-iv19-complex-hom-lattice}
\notready
\end{remark}

\begin{remark}[Rosati trace package]
\label{mumford-iv21-rosati-positive-trace}
The reduced-trace pairing $(u,v)\mapsto\Trd(u^\dagger v)$ is positive
definite on the real endomorphism algebra.
\dcref{M:ch4:IV.21}
\uses{mumford-thm-rosati-positive}
\notready
\end{remark}

\begin{remark}[Theta-group definition]
\label{mumford-iv23-theta-group-definition}
The theta group is the central extension of $K(L)$ by scalar automorphisms of
the line bundle $L$.
\dcref{M:ch4:IV.23}
\uses{mumford-def-theta-group,mumford-iv23-line-bundle-theta-group}
\notready
\end{remark}

\section{Riemann forms}
\label{mumford-IV20}
\dcref{M:ch4:IV.20}

\begin{definition}[The $\ell$-adic Weil pairing]
\label{mumford-def-weil-pairing-ladic}
\dcref{M:ch4:IV.20}
For $\ell\ne\operatorname{char}(k)$, the Poincare bundle induces compatible
perfect pairings
\[
 e_{\ell^n}:A[\ell^n]\times\Adual[\ell^n]\to\mu_{\ell^n}.
\]
Their inverse limit is a perfect $\Zell$-bilinear pairing
$T_\ell(A)\times T_\ell(\Adual)\to\Zell(1)$.
\uses{mumford-prop-weil-pairing,mumford-def-tate-module}
\notready
\end{definition}

\begin{definition}[Riemann form of a polarization]
\label{mumford-def-riemann-form-algebraic}
\dcref{M:ch4:IV.20}
For a polarization $\lambda:A\to\Adual$, the finite-level multiplicative
pairings are
\[
 e_{\ell^n}^{\lambda}(x,y)=e_{\ell^n}(x,\lambda(y))
\]
on $A[\ell^n]\times A[\ell^n]$; their inverse limit is a pairing with values
in $\Zell(1)$.  After choosing a compatible generator of the Tate twist (or
passing to the associated additive lattice), this gives an alternating
$\Zell$-valued form.  It is nondegenerate over $\Qell$ and positive exactly
when $\lambda$ comes from an ample line bundle.
\uses{mumford-def-polarization,mumford-def-weil-pairing-ladic}
\notready
\end{definition}

\begin{theorem}[Homomorphisms represented by Riemann forms]
\label{mumford-thm-riemann-form-hom}
For $u:A\to\Adual$, the bilinear form
$(x,y)\mapsto e_\ell(x,u(y))$ is skew-symmetric if and only if
$u=\phi_L$ for a line bundle $L$ on $A$.  If $L$ is ample, this form is a
polarization Riemann form.
\dcref{M:ch4:IV.20}
\uses{mumford-prop-phi-L,mumford-def-weil-pairing-ladic}
\notready
\end{theorem}

\section[Positivity of Rosati]{Positivity of the Rosati involution}
\label{mumford-IV21}
\dcref{M:ch4:IV.21}

\begin{definition}[Rosati involution]
\label{mumford-def-rosati}
\dcref{M:ch4:IV.21}
For a polarization $\lambda:A\to\Adual$, the Rosati involution on
$\End^0(A)$ is
\[
 u^\dagger=\lambda^{-1}\circ\widehat u\circ\lambda.
\]
It is a positive involution: $u^\dagger u$ has positive reduced trace for
$u\ne0$.
\uses{mumford-def-polarization,mumford-prop-biduality-any-characteristic}
\notready
\end{definition}

\begin{theorem}[Trace-form positivity]
\label{mumford-thm-rosati-positive}
Using the reduced trace on each simple factor, the bilinear form
\[
 (u,v)\longmapsto\Trd_{\End^0(A)/\Q}(u^\dagger v)
\]
is symmetric and positive definite on $\End^0(A)\otimes_\Q\R$.  The
Rosati-symmetric subspace identifies with $\NS(A)\otimes\Q$, and ample classes
are exactly its positive elements.
\dcref{M:ch4:IV.21}
\uses{mumford-def-rosati,mumford-thm-riemann-form-hom,mumford-prop-pic-ns-sequence}
\notready
\end{theorem}

\begin{corollary}[Semisimplicity up to isogeny]
\label{mumford-cor-semisimple-endomorphism}
\dcref{M:ch4:IV.21}
The algebra $\End^0(A)$ is semisimple, and its simple factors are division
algebras equipped with positive involutions induced by Rosati.  The isogeny
decomposition of $A$ is unique up to permutation and isogeny of factors.
\uses{mumford-thm-rosati-positive,mumford-thm-poincare-complete-reducibility}
\notready
\end{corollary}

\section{Examples}
\label{mumford-IV22}
\dcref{M:ch4:IV.22}

\begin{example}[Elliptic curves]
\label{mumford-ex-elliptic-endomorphisms}
\dcref{M:ch4:IV.22}
For an elliptic curve $E$, $\End^0(E)$ is either $\Q$, an imaginary quadratic
field (complex multiplication), or a quaternion division algebra in the
supersingular characteristic-$p$ case.  The Rosati involution is complex
conjugation in the middle case and the standard positive involution in the
quaternionic case.
\uses{mumford-cor-simple-endomorphism-algebra,mumford-def-rosati}
\notready
\end{example}

\begin{example}[Products and matrix algebras]
\label{mumford-ex-product-endomorphisms}
\dcref{M:ch4:IV.22}
If $A=B^r$ with $B$ simple, then
$\End^0(A)\simeq M_r(\End^0(B))$ and the product polarization induces the
adjoint-transpose Rosati involution.  Non-isogenous simple factors give a
product decomposition of the endomorphism algebra.
\uses{mumford-cor-semisimple-endomorphism,mumford-thm-poincare-complete-reducibility}
\notready
\end{example}

\section[Theta groups]{The theta group of a line bundle}
\label{mumford-IV23}
\dcref{M:ch4:IV.23}

\begin{definition}[Theta group]
\label{mumford-def-theta-group}
\dcref{M:ch4:IV.23}
For a line bundle $L$ on $A$, the theta group $\mathcal G(L)$ is the group
scheme whose $S$-points are pairs $(x,\varphi)$ with $x\in K(L)(S)$ and
$\varphi:L_S\xrightarrow{\sim}t_x^*L_S$.  There is an exact sequence
\[
 1\longrightarrow\Gm\longrightarrow\mathcal G(L)\longrightarrow K(L)
 \longrightarrow1.
\]
\uses{mumford-def-KL}
\notready
\end{definition}

\begin{proposition}[Theta commutator]
\label{mumford-prop-theta-commutator}
\dcref{M:ch4:IV.23}
The commutator in the theta group is the alternating pairing
\[
 e^L:K(L)\times K(L)\to\Gm
\]
induced by the biextension of $L$.  If $L$ is nondegenerate, this pairing is
perfect and $|K(L)|$ is a square.
\uses{mumford-def-theta-group,mumford-def-biextension}
\notready
\end{proposition}

\begin{theorem}[Theta representation]
\label{mumford-thm-theta-representation}
When $L$ is ample, $\mathcal G(L)$ acts on $H^0(A,L)$.  Every irreducible
representation on which the central $\Gm$ acts by its tautological character
has dimension $h^0(A,L)$ and is unique up to isomorphism; the natural action on
$H^0(A,L)$ is that representation.
\dcref{M:ch4:IV.23}
\uses{mumford-prop-theta-commutator,mumford-thm-line-bundle-index}
\notready
\end{theorem}

\section[Complex comparison]{The complex case}
\label{mumford-IV24}
\dcref{M:ch4:IV.24}

\begin{theorem}[Analytic and algebraic Riemann forms agree]
\label{mumford-thm-riemann-comparison}
For a complex abelian variety $A(\C)=V/\Lambda$ and an ample line bundle $L$,
the algebraic Riemann form obtained from the Weil pairing is the integral
alternating form $\operatorname{Im}H$ in the Appell--Humbert description of the
analytic bundle.  Under this identification, the algebraic Rosati involution
is the adjoint involution for the positive Hermitian form $H$.
\dcref{M:ch4:IV.24}
\uses{mumford-thm-riemann-form-hom,mumford-thm-rosati-positive,mumford-thm-appell-humbert}
\notready
\end{theorem}

\begin{corollary}[Complex polarization dictionary]
\label{mumford-cor-complex-polarization-dictionary}
\dcref{M:ch4:IV.24}
Polarizations of $A$ correspond to positive integral Riemann forms on
$(V,\Lambda)$, principal polarizations to unimodular forms, and the theta
group of $L$ is the analytic Heisenberg extension attached to the finite
symplectic module $K(L)$.
\uses{mumford-thm-riemann-comparison,mumford-thm-theta-representation}
\notready
\end{corollary}
